\documentclass[11pt]{article}
\usepackage[margin=1.05in]{geometry}
\usepackage{amsmath,amssymb,amsthm}
\usepackage[hidelinks]{hyperref}
\newtheorem{theorem}{Theorem}
\newcommand{\bin}[2]{\binom{#1}{#2}}
\title{Splitting a bound can destroy interlacing\protect\\for slices of prisms}
\author{}
\date{25 November 2026}
\begin{document}
\maketitle
\begin{abstract}
Ferroni and McGinnis conjectured that splitting a coordinate bound of
an integral prism slice produces interlacing Ehrhart $h^*$-polynomials.
We give counterexamples for every central slice of a box with bounds
$(1^{2a-2},2)$ and $a\ge6$. For $a\ge7$, Ehrhart reciprocity and
Vieta's formula yield a contradiction using only the top two coefficients.
At $a=6$ both polynomials have simple negative roots, and an integer
separates their leftmost roots in the forbidden order.
\end{abstract}

For positive integers $c_1,\ldots,c_s$ and $0<k<\sum_i c_i$, put
\[
 R_{k,c}=\{x\in\mathbb R^s:0\le x_i\le c_i,\quad \textstyle\sum_i x_i=k\}.
\]
This is a lattice polytope of dimension $D=s-1$. Write $H_{k,c}$ for its
Ehrhart $h^*$-polynomial, so that its lattice counts $L(t)$ satisfy
$\sum_{t\ge0} L(t)z^t=H_{k,c}(z)/(1-z)^{D+1}$.
Ferroni and McGinnis \cite[Conjecture~5.12]{fm} ask whether all these
numerators are real-rooted, and whether splitting $c_s\ge2$ into
$c_s-1,1$ makes the two numerators interlace. We disprove the latter
assertion. The individual real-rootedness question remains open.

\begin{theorem}
For each integer $a\ge6$, the polynomials
$P=H_{a,(1^{2a-2},2)}$ and $Q=H_{a,(1^{2a})}$ do not interlace.
\end{theorem}

\begin{proof}
The two slices are central. Since their centers are interior to their
boxes, the relative interior of a slice is given by strict coordinate
bounds. The dimensions of the two slices are $D_P=2a-2$ and $D_Q=2a-1$.
Let $L^\circ$ count the lattice points in the relative interior of a
positive integer dilation. Ehrhart reciprocity gives
\begin{equation}\label{recip}
 \sum_{t\ge1}L^\circ(t)z^t=
   \frac{z^{D+1}H_{k,c}(1/z)}{(1-z)^{D+1}}.
\end{equation}
For both slices $L^\circ(1)=0$ and $L^\circ(2)=1$. In the second
dilation all unit-bound coordinates must be one, and the sum forces the
exceptional coordinate of the first slice to be two. By
\eqref{recip}, each numerator is monic of degree $D-1$, and its next
leading coefficient is $L^\circ(3)-(D+1)$.

Put $b=a-1$. In the third dilation of the first slice the exceptional
coordinate ranges from one to five. The other $2b$ coordinates are one
or two and their sum determines the exceptional coordinate. Therefore
\[
 L_P^\circ(3)=\bin{2b}{b}+2\bin{2b}{b-1}+2\bin{2b}{b-2}.
\]
In the second slice exactly $a$ of the $2a$ coordinates are two, so
$L_Q^\circ(3)=\bin{2a}{a}$. Write $p$ and $q$ for the next leading
coefficients of $P$ and $Q$. Subtracting, we obtain
\begin{equation}\label{gap}
 p-q=1+\bin{2b}{b}\frac{b^2-5b-2}{(b+1)(b+2)}.
\end{equation}
This is positive for $b\ge6$.

If the roots interlaced, they would all be negative, by nonnegativity
of the Ehrhart numerators and their positive constant terms. Since
$Q$ has one higher degree, writing the zeros in increasing order
would give $v_1\le u_1\le v_2\le\cdots\le u_m\le v_{m+1}<0$,
with the $u_i$ from $P$. Hence
\[
 p=-\sum_{i=1}^m u_i\ \le\ -\sum_{i=1}^m v_i
 \ <\ -\sum_{i=1}^{m+1}v_i=q.
\]
This contradicts \eqref{gap} for $a\ge7$. Already at $a=7$ the two
coefficients are $3485$ and $3418$.

At $a=6$ the Ehrhart numerators have the following coefficients.
\begin{center}\small
\begin{tabular}{r|r|r@{\qquad}r|r|r}
$i$&$[z^i]P$&$[z^i]Q$&$i$&$[z^i]P$&$[z^i]Q$\\\hline
0&1&1&6&347986&3834832\\
1&661&912&7&38686&878932\\
2&31846&62767&8&901&62767\\
3&312146&878932&9&1&912\\
4&927706&3834832&10&0&1\\
5&960774&6169360&&&\end{tabular}
\end{center}
For explicit reproduction, the entries are
\[
 [z^j]H_{k,c}(z)=\sum_{t=0}^j(-1)^{j-t}\bin{s}{j-t}
       [u^{kt}]\prod_{i=1}^s(1+u+\cdots+u^{tc_i}).
\]
We have $P(-850)=1625988474636775883423151>0$. Its odd degree and
positive leading coefficient force a root below $-850$. For the other
numerator, direct binomial expansion of the displayed table gives
\[
 [y^i]Q(-850-y)=(-1)^i\sum_{j=i}^{10}\bin ji(-850)^{j-i}[z^j]Q(z)>0
 \quad(0\le i\le10).
\]
Thus $Q$ has no root at or below $-850$. The smaller-degree polynomial
has a root to the left of every real root of the larger-degree one,
which rules out interlacing.
\end{proof}

\paragraph{Exact checks.}
The accompanying certificate gives disjoint rational sign-change intervals
for all nine and all ten roots at $a=6$. The two leftmost intervals are
$(-857,-856)$ and $(-839,-838)$. Each polynomial separately has only
simple negative roots. The verifier obtains all lattice counts again by
bounded-composition inclusion--exclusion and evaluates endpoint signs
using integer or rational arithmetic. The argument for $a\ge7$ requires
no computation of roots.

\begin{thebibliography}{9}
\bibitem{fm} L. Ferroni and D. McGinnis, Lattice points in slices of prisms,
\emph{Canadian Journal of Mathematics} 77 (2025), 1013--1040.
\href{https://doi.org/10.4153/S0008414X24000233}{doi:10.4153/S0008414X24000233}.
\end{thebibliography}
\end{document}
